\documentclass[11pt,a4paper]{article}
\usepackage[utf8]{inputenc}
\usepackage[T1]{fontenc}
\usepackage[italian]{babel}
\usepackage{lmodern}
\usepackage[margin=1.6cm,top=1.3cm,bottom=1.5cm]{geometry}
\usepackage{amsmath,amssymb,tikz}
\usepackage{fancyhdr}
\setlength{\parindent}{0pt}
\setlength{\parskip}{3pt}
\setlength{\abovedisplayskip}{5pt}
\setlength{\belowdisplayskip}{5pt}
\setlength{\emergencystretch}{2em}
\pagestyle{fancy}\fancyhf{}
\fancyfoot[L]{\small 1SP -- Razionali e percentuali}
\fancyfoot[R]{\small\thepage/2}
\renewcommand{\headrulewidth}{0pt}
\renewcommand{\footrulewidth}{0pt}
\newcommand{\blocco}[3]{\par\medskip\noindent\textbf{#1. #2}\hfill{\small #3 min}\par\nopagebreak}
\newcommand{\terna}[3]{%
\par\smallskip\noindent
\begin{minipage}[t]{.323\textwidth}\raggedright
\textbf{\small A -- Insieme alla lavagna}\par\smallskip #1
\end{minipage}\hfill
\begin{minipage}[t]{.323\textwidth}\raggedright
\textbf{\small B -- Subito da solo}\par\smallskip #2
\end{minipage}\hfill
\begin{minipage}[t]{.323\textwidth}\raggedright
\textbf{\small C -- Da solo}\par\smallskip #3
\end{minipage}\par}
\newcommand{\retta}{%
\par\noindent\hfill
\begin{tikzpicture}[x=4mm,y=4mm]
\path[use as bounding box] (-13,-2) rectangle (25,2);
\draw[step=1,gray!35,very thin] (-12,-1) grid (24,1);
\draw[->,line width=.65pt] (-12.6,0)--(24.6,0);
\foreach \x in {-12,...,24}{\draw[line width=.4pt] (\x,-.14)--(\x,.14);}
\foreach \x/\n in {-12/-1,0/0,12/1,24/2}{
\draw[line width=.7pt] (\x,-.3)--(\x,.3);
\node[below,font=\small] at (\x,-1.05) {$\n$};}
\end{tikzpicture}\hfill\mbox{}\par}
\begin{document}

\begin{center}
{\Large\textbf{Razionali e percentuali}}\\[2pt]
\textbf{1SP -- Esercizi guidati e autonomi -- 50 minuti}
\end{center}
\textbf{Nome e cognome:} \dotfill \quad \textbf{Data:} \dotfill

\textbf{Come lavoriamo.} Per ogni blocco facciamo A insieme; subito dopo svolgi B e C da solo. Scrivi i calcoli sul quaderno e usa le rette della scheda. Senza calcolatrice.

\blocco{1}{Somma algebrica: frazioni e parentesi}{8}
Calcola e riduci ai minimi termini. Risolvi prima le parentesi interne e controlla i cambi di segno.
\terna{
\[
\begin{aligned}
&\frac56-\left[\frac34-\left(\frac23-\frac12\right)\right]\\[4pt]
&\qquad+\frac15-\frac7{10}
\end{aligned}
\]
}{
\[
\begin{aligned}
&\frac76-\left[\frac54-\left(\frac23-\frac12\right)\right]\\[4pt]
&\qquad+\frac25-\frac3{10}
\end{aligned}
\]
}{
\[
\begin{aligned}
&\frac{11}{12}-\left[\frac78-\left(\frac56-\frac34\right)\right]\\[4pt]
&\qquad+\frac35-\frac1{10}
\end{aligned}
\]
}

\blocco{2}{Esponenti negativi e proprietà delle potenze}{9}
Riduci prima prodotti, quozienti e potenze di potenze. Calcola poi gli inversi e la somma algebrica finale.
\[
\begin{aligned}
\text{\small\bfseries A -- Insieme:}\quad&
\left[\left(-\frac23\right)^{-2}\cdot\left(-\frac23\right)^3\right]
:\left(-\frac23\right)^2+\left(\frac52\right)^{-1}\\[8pt]
\text{\small\bfseries B -- Da solo:}\quad&
\left[\left(-\frac23\right)^4:\left(-\frac23\right)^6\right]
\cdot\left(-\frac23\right)^3-\left(\frac52\right)^{-1}\\[8pt]
\text{\small\bfseries C -- Da solo:}\quad&
\left[\left(-\frac23\right)^{-1}\right]^2
:\left(-\frac23\right)^{-3}+\left(\frac53\right)^{-1}
\end{aligned}
\]

\blocco{3}{Ordina e rappresenta sulla retta}{8}
Ordina sul quaderno usando $<$ oppure $=$, poi rappresenta ogni numero sulla retta: frazioni equivalenti occupano lo stesso punto. Il m.c.m. dei denominatori è \textbf{12}: \textbf{1 unità = 12 quadretti}, un quadretto vale $1/12$.

\textbf{A -- Insieme:}\quad
$\displaystyle -\frac7{12},\quad\frac56,\quad-\frac34,\quad\frac76,\quad\frac23,\quad\frac8{12}$
\retta
\textbf{B -- Da solo:}\quad
$\displaystyle -\frac56,\quad\frac7{12},\quad-\frac23,\quad\frac54,\quad\frac36,\quad\frac24$
\retta
\textbf{C -- Da solo:}\quad
$\displaystyle -\frac{11}{12},\quad\frac76,\quad-\frac34,\quad\frac9{12},\quad\frac53,\quad\frac34$
\retta

\newpage

\begin{center}
{\large\textbf{Dalla virgola alle percentuali}}
\end{center}

\blocco{4}{Dalla virgola alla frazione}{6}
Trasforma in frazioni ridotte ai minimi termini. La barretta indica le cifre che si ripetono.
\terna{
\[
0{,}2 \qquad 0{,}\overline3
\]
}{
\[
0{,}4 \qquad 0{,}\overline6
\]
}{
\[
0{,}6 \quad 0{,}8 \quad 0{,}\overline2
\]
}

\medskip
\textbf{L'idea pratica.} $p\%$ significa $p$ su 100, cioè $\dfrac{p}{100}$.
Per trovare una parte, moltiplica il totale per questa frazione; per trovare la percentuale, dividi la parte per il totale e moltiplica per 100.

\blocco{5}{Frazioni note, decimali e percentuali}{4}
Completa ogni catena. Cerca rappresentazioni equivalenti con denominatore 100.
\terna{
\[
\frac15=\ldots=\ldots\%
\]
\[
\frac12=\ldots=\ldots\%
\]
}{
\[
\frac25=\ldots=\ldots\%
\]
\[
\frac14=\ldots=\ldots\%
\]
}{
\[
\frac35=\ldots=\ldots\%
\]
\[
\frac45=\ldots=\ldots\%
\]
\[
\frac34=\ldots=\ldots\%
\]
}

\vspace{4mm}
\blocco{6}{Calcolare una percentuale di un totale}{5}
Scrivi la percentuale come frazione e calcola la parte richiesta.
\terna{
In un allenamento una squadra effettua \textbf{50 tiri} e ne realizza il \textbf{20\%}. Quanti canestri segna?
}{
In un allenamento una squadra effettua \textbf{50 tiri} e ne realizza il \textbf{40\%}. Quanti canestri segna?
}{
In un allenamento una squadra effettua \textbf{50 tiri} e ne realizza il \textbf{60\%}. Quanti canestri segna?
}
\par\medskip
\textbf{Controllo:} la percentuale si riferisce a quale totale? La parte trovata è plausibile?

\vspace{4mm}
\blocco{7}{Trovare la percentuale conoscendo parte e totale}{4}
Scrivi prima il rapporto \emph{canestri/tiri}, riducilo e trasformalo in percentuale.
\terna{
Una giocatrice segna \textbf{8 canestri su 20 tiri}. Qual è la percentuale di realizzazione?
}{
Una giocatrice segna \textbf{12 canestri su 20 tiri}. Qual è la percentuale di realizzazione?
}{
Una giocatrice segna \textbf{16 canestri su 20 tiri}. Qual è la percentuale di realizzazione?
}

\vspace{4mm}
\blocco{8}{Uno sconto: dalla percentuale al prezzo finale}{4}
Calcola prima lo sconto in euro, poi il prezzo da pagare.
\terna{
Una maglia costa \textbf{50 euro}. Lo sconto è del \textbf{20\%}. Quanto risparmi? Quanto paghi?
}{
Una maglia costa \textbf{50 euro}. Lo sconto è del \textbf{40\%}. Quanto risparmi? Quanto paghi?
}{
Una maglia costa \textbf{50 euro}. Lo sconto è del \textbf{60\%}. Quanto risparmi? Quanto paghi?
}

\vfill
\textbf{Numeri da riconoscere al volo.} Ricostruisci senza calcolatrice:
\[
\frac15,\quad\frac25,\quad\frac35,\quad\frac45
\quad\longleftrightarrow\quad
\text{decimale finito}\quad\longleftrightarrow\quad\text{percentuale};
\qquad
0{,}\overline3=\frac13,\quad 0{,}\overline6=\frac23.
\]
\textbf{Controlla i quinti:} $3/5=0{,}6$ e $4/5=0{,}8$.

\end{document}
